% GENERATED FILE -- do not edit by hand.
% Source: tests/test_defence_models.py (IR model definitions)
% SHA-256: test-code
% Run: python scripts/generate_paper_figures.py

\begin{figure}[t]
\centering
\begin{tikzpicture}[
  node distance=1.5cm and 2cm,
  state/.style={draw, circle, minimum size=1cm, font=\footnotesize},
  bad/.style={fill=unsafebg!60, draw=unsafebd},
  good/.style={fill=safebg!60, draw=safebd},
  model/.style={draw, rounded corners, minimum width=3.5cm, minimum height=0.8cm,
               font=\small, align=center},
  arrow/.style={-Stealth, thick},
]
% Colors defined globally in preamble

% --- Dual-LLM ---
\node[model] (dllm) at (0, 2) {Dual-LLM baseline};
\node[model, good] (dllm_q) at (-2, 0.5) {Property Q:\\processor never executes\\{\bf SAFE}};
\node[model, bad] (dllm_pe) at (2, 0.5) {PE property:\\attacker influences action\\{\bf UNSAFE}};

% --- Requester-only ---
\node[model] (req) at (7, 2) {Requester-only\\(defective)};
\node[model, bad] (req_pe) at (7, 0.5) {PE property:\\ignores attacker\\{\bf UNSAFE}};

% --- ITES reference ---
\node[model] (ites) at (12, 2) {ITES reference};
\node[model, good] (ites_pe) at (12, 0.5) {PE property:\\all principals checked\\{\bf SAFE}};

% Arrows
\draw[arrow] (dllm) -- (dllm_q);
\draw[arrow] (dllm) -- (dllm_pe);
\draw[arrow] (req) -- (req_pe);
\draw[arrow] (ites) -- (ites_pe);

% Annotation
\node[font=\footnotesize, align=center] at (0, -1) {Satisfying Q does not\\imply PE safety};
\node[font=\footnotesize, align=center] at (12, -1) {Intersection rule\\preserves PE};
\end{tikzpicture}
\caption{Comparative defence verification on finite IR models.  The
Dual-LLM baseline satisfies its own property Q (processor never
executes) but violates the privilege-escalation property.  The
requester-only negative control also violates PE.  The ITES reference
preserves PE.  All results are finite IR models, not
implementation-conformance evidence.}
\label{fig:defence-comparison}
\end{figure}
